% SAME-FAMILY GENERAL COMPACT SOURCE DRAFT, 25 October2026.
% Exact19 whole A/B/D consumers; no new paper, canonical or main edit.
% No arithmeticity, field/integrality, enumeration, imports, tests or builds.
% DRAFT pending complete different-author/root review and shared-tail closure.
\section{The whole classical forests at exact dimension nineteen}
\label{app:exact19:wholeclassical}

These are three actual-rank consumers for
the general compact argument.
Every point and pressure below belongs
to one original farthest vertex \(p\),
chosen from one strict interior point.
We keep its single decomposition
\(z=E(p-x)\), \(x\in p^\perp\),
\(\|x\|<1\).
No arithmetic restriction is used.

\subsection{A whole \(A_{19}\): the window and the three reference shapes}
\label{app:exact19:Awindow}

Use the zero-sum subspace of
\(\mathbb R^{20}\), with actual old roots
and duals
\[
 n_i=(e_i-e_{i+1})/\sqrt2,\quad
 \Omega_i=\sqrt2(P_i-iP_{20}/20),\quad
 D_i=\Omega_i^2=i(20-i)/10 .
\]
Write \(h_i=x\cdot\Omega_i\), \(h_0=h_{20}=0\).
The sequence is strictly concave and
\[
                 \|x\|^2=\frac12\sum_{i=0}^{19}(h_{i+1}-h_i)^2 .
\]
A central site \(5\le i\le15\) is weak
when \(h_i\le\sqrt{D_i}-.61\), equality
included.

We use the shared complete \(.61\) window
of Lemma~\ref{app:tail:lem:Awindow2021}
at \(N=20\), before any contrary
component-energy bound: every weak cap
is pure, with no spectator projection,
and is either
\[
 v_i=g\Omega_i+\Omega_e/2,\quad e=1\text{ or }19,
 \quad \cos(\pi/5)\le g<27/25 ,
\]
or a near-four cap at \(i=5,6\) or
a reflected site.
The first row has own native five,
finite native at least seven, or dotted;
the smaller own labels and parallel case
have their actual affine or feasible-null
obstructions.
The literal helpers are nonadjacent.
These are the actual finite-white,
own-zero, foreign-clearing and mixed
two-cap conclusions of the shared
window proof, with retained paths of
rank at least four and full connected
rank nineteen; they do not follow by
applying a higher-rank whole-forest
exclusion here.
Its full internal proof precedes
the higher-tail energy consumers;
it is not a published classification
or an additional assumed field condition.

There is a useful stronger removal of
near four at this rank.
For a left near-four cap at \(i=5,6\),
its positive finite replacement at
\(j=i-1\) is an actual endpoint.
Its white forest is
\(B_{21-i}+A_{i-2}\); positive inverse
projection proves feasibility against
every original wall, including the
released \(j\).
The exact edge cosine is
\[
 \cosh L_j=\sqrt{2(i-1)/20}
                 \left\{(21-i)/2+(20-i)/\sqrt2\right\}>11 .
\]
For \(i=5\), use
\(.632(8+7.5(1.414))>11\);
for \(i=6\), use
\(.707(7.5+7(1.414))>11\).
Its normalized dual demand is
\(>\sqrt{5/6}>.91\).
At the opposite end, either a second
near-four mate has that same demand,
or helper nonadjacency gives a strong
\(k=14\) or fifteen, with demand
\(>1-.61/2.7>.774\).
The correct normalized correlation is
\[
 \rho=\sqrt{\frac{j(20-k)}{(20-j)k}}
                       \le1/\sqrt7<.38 .
\]
Cauchy would give
\((.91+.774)^2<2+2(.38)=2.76\);
instead \(1.684^2=2.835856\).
Thus near four is impossible under
\(\|x\|<1\), without a stronger
component assumption.

All remaining weak sites are therefore
nonadjacent pure helpers.
Their edge-ratio target is
\[
 R(q)=100/61-\frac{q-.61}{q(2q-.61)} .
\]
It increases for \(q\ge5/2\), while
\(q\le\sqrt{10}<16/5\) and
\(R(16/5)=3/2-251/565104<3/2\).
For side order \(m\ge5\),
\(\sigma=g+1/(2m)\), \(\delta=(m+1)/(2m)\),
and
\[
 \frac{\sigma^2}{\delta}-(4g-2)
                   =\frac{(g-1-1/(2m))^2}{\delta}\ge0 .
\]
Weakness consequently sharpens the
entire helper range to \(g<17/16\).
Native five has \(.809<g<.81\);
every other finite own has \(.9<g<1\).

Take rational symmetric heights
\[
 H_5=2.128,\ H_6=2.288,\ H_7=2.406,\
 H_8=2.488,\ H_9=2.536,\ H_{10}=2.552,\qquad
 H_{20-i}=H_i .
\]
They lie strictly below the strong
floors by positive rational squaring.
Linear interpolation on the tails
gives nonnegative dual coefficients.
Omit every actual internal weak knot
by replacing its height by the mean
of its neighbours.
The holes are nonadjacent, so their
open affected edge pairs are disjoint
and their exact quarter-square losses
add.
If an end site five or fifteen is
weak, flatten that entire end through
six or fourteen instead.
The complete reference data are
\[
\begin{array}{c|c|c|c}
 \text{weak end count}&E=\|w\|^2\text{ at least}
          &\text{end block lengths}&\text{individual capacities}\\ \hline
 0&.9544848-5(.000441)=.9522798&5,5&
             \beta_5,\beta_{15}\ge83/625\\
 1&3462479/3750000&6,5&
             \beta_6\ge79/600,\ \beta_{15}\ge83/625\\
 2&1341604/1500000&6,6&
             \beta_6,\beta_{14}\ge79/600 .
\end{array}
\tag{A19refs}\label{app:exact19:eq:Arefs}
\]
For count one orient the weak end to
the left by relabelling the old path.
This assumes no symmetry of the polytope.
There are at most five internal holes
in the first row, four in the other
rows.
The one-end row loses at most
\(4(.000441)\); the two-end row loses
at most \(4(.000324)\).
Their exact undeleted energies follow
from
\[
 a^2/5+(b-a)^2+T,\quad
 \tfrac12[b^2/6+a^2/5+(b-a)^2+2T],\quad
 b^2/6+T,
\]
where \(a=2.128\), \(b=2.288\) and
\(T=.118^2+.082^2+.048^2+.016^2\).
The boundary coefficients are
\[
 (a/5-(b-a))/2=83/625,\qquad
                    ((7/6)b-H_7)/2=79/600 .
\]
An adjacent hole increases these
capacities.
All positive coefficients are paid
by actual strong floors, so
\(x\cdot w\ge E\) in every allowed
hole pattern.
The constant signed means obey
\[
                  .3<b_5<.302,\qquad1/4<b_6<.27 .
\]

\subsection{The actual helper feet and their complete continuous bounds}
\label{app:exact19:Ahelperfeet}

For a weak site five, write
\(c=Ap-v\), \(v=g\Omega_5+\Omega_e/2\).
The fresh exact times are
\[
 A_{\rm far}^2=(300g^2+20g-21)/40,\qquad
 A_{\rm near}^2=(300g^2+60g-21)/40 .
\tag{A19times}\label{app:exact19:eq:Atimes}
\]
Both formulas concern \(N=20\).

The following constructions give actual
pure spatial pressures \(x\cdot U\ge H\).
At native five use the finite \(H_4\)
correction
\[
 k=g n_5+\tfrac23n_e+\tfrac13n_{\rm adj},\qquad
 \delta=2/3-g^2>0,\quad U=v-k .
\]
The normalized finite projected point
is obtained from \(p+(A/\delta)(c+k)\).
Its future time is positive, its square
is \(-1-A^2/\delta\), and all selected
products vanish.
Every omitted original-wall product
is nonpositive by the positive
normal coefficients and acuteness.
The pressure is
\[
 H=A+\delta/A-\sqrt{\delta(A^2+\delta)}/A
                   >A-\sqrt\delta>A-.111 .
\]
The spacelike correction \(c+k\)
alone is not a hyperbolic point.

For finite own at least seven,
put \(d=1-g^2\) and use
\(\xi=c+g n_5+2d n_e\).
Its selected products are zero,
zero and \(-1/2+2d<0\);
its square is \(-d(1-4d)<0\).
It is future and satisfies every
other original wall termwise.
Thus \(U=v-g n_5-2d n_e\),
\(H=A-\sqrt{d(1-4d)}\), with the
last radical at most \(1/4\).
For dotted own use
\(\xi=c+g^{-1}n_5\), of square
\(-(1-g^{-2})\), zero cap product
and negative own-root product.
It gives \(U=v-g^{-1}n_5\),
\(H=A-\sqrt{1-g^{-2}}\).
All vectors \(U\) are spatial in
the original \(p^\perp\).

A near helper forces \(x_1>.64\).
Indeed weakness gives \(h_5<2.129\),
and dropping the positive old-root
corrections from the pressure yields
\(x_1\ge\sqrt2(H-gh_5)\).
At native five, \(A>2.365\), giving
\[
              x_1>1.414(2.365-.111-.81(2.129))>.64 .
\]
For finite \(g\ge.9\),
\(A_{\rm near}'(g)>\sqrt{7.5}>2.7\);
thus \(A_{\rm near}-2.129g\) increases.
At \(.9\), \(A>2.626\), giving
\(1.414(2.626-.25-.9(2.129))=.6502986\).
For dotted \(1<g<17/16\),
\(A_{\rm near}'<3\), whereas
\[
                  \frac{d}{dg}\sqrt{1-g^{-2}}
                  =\frac1{g^3\sqrt{1-g^{-2}}}>2 .
\]
Here \(g^3<5/4\), the radical \(<2/5\).
The difference \(A_{\rm near}-\sqrt{1-g^{-2}}-2.129g\)
decreases.
At \(17/16\), its time is \(>3.0879\)
and radical \(<.338\), which again
gives \(x_1>.64\).
All time/derivative comparisons
follow by positive squaring in
\eqref{app:exact19:eq:Atimes}.

If site five is weak, its first-six
reference block has mean \(b_6<.27\).
Grafting \(.64\) into its first
coordinate, preserving the sum,
raises its energy by
\((6/5)(.64-.27)^2=.16428\).
The other five coordinates remain
\(>(1.5-.64)/5=.172\), above the
next old coordinate \(<.118\).
This is a legal ordered dual
reference, and even the weakest
row of \eqref{app:exact19:eq:Arefs}
then exceeds one.
Every weak end helper is therefore
FAR, in both orientations.

For a far helper on six/six blocks,
block averaging gives
\[
 \Pi U=(5g/6)\Omega_6+\Omega_{14}/12,\quad
 Y=U-\Pi U,\quad w\cdot Y=0,\quad
              \Delta=H-2.288(5g/6+1/12).
\]
The three residual squares are
\[
\begin{array}{c|c}
 \text{own}&Q=\|Y\|^2\\ \hline
 5&(2/3)g^2+1/12\\
 \text{finite }\ge7&4g^4-(16/3)g^2+29/12\\
 \text{dotted}&(5/3)g^2-19/12+g^{-2}.
\end{array}
\tag{A19Q}\label{app:exact19:eq:AhelperQ}
\]
At native five, \(A_{\rm far}>2.188\),
so \(\Delta>.34>.30\), \(Q<.521<13/12\).
For all other finite owns,
\(A_{\rm far}-2.288(5g/6+1/12)\)
increases, since
\(A_{\rm far}'>\sqrt{7.5}>5(2.288)/6\).
The complete radical/time seals are
\[
\begin{array}{c|c|c|c}
 g\text{ interval}&A_{\rm far}\text{ at left}>
      &\sqrt{(1-g^2)(4g^2-3)}\le&\Delta>\\ \hline
 (.9,.91]&2.449&.2318&.3105333\ldots\\
 [.91,.92]&2.478&.2434&.3088666\ldots\\
 [.92,1)&2.5065&.25&.3117 .
\end{array}
\]
The first two radical bounds use
the increasing quartic on those
short intervals; its global maximum
is \(1/16\).
The tight middle lower bound is
\(23165/75000\).
The finite residual polynomial
increases for \(g^2>.81\) and
tends to \(13/12\).
Thus every finite case has
\(\Delta>.30\), \(Q<13/12\).

In the dotted interval the same
derivative argument makes \(\Delta\)
decrease.
At \(17/16\), \(A_{\rm far}>2.9108\),
the radical is \(<.338\), and
\(2.288(5g/6+1/12)=2.2165\).
Thus \(\Delta>.3563>.35\).
Its \(Q\) increases with \(g^2\ge1\)
and at \(17/16\) equals
\(229/768+256/289<148/125\);
the cleared margin is \(271>0\).
These bounds cover the whole
continuous dotted interval.

\subsection{Two weak ends: the exact nonzero cross and both prices}
\label{app:exact19:Atwoweak}

Suppose both five and fifteen
are weak.
Use the six/six reference in
\eqref{app:exact19:eq:Arefs}.
For a left far-helper foot,
its scaled first and last residual
blocks are
\[
\begin{split}
 \sqrt2Y_{\rm own}&=(g/3,g/3,g/3,g/3,g/3-k,-5g/3+k),\\
 \sqrt2Y_{\rm helper}&=(1/6,1/6,1/6,1/6-m,1/6-l+m,-5/6+l).
\end{split}
\]
Here \((k,l,m)\) equals
\((g,2/3,1/3)\) at native five,
\((g,2(1-g^2),0)\) for other finite
owns, and \((1/g,0,0)\) for dotted.
The blocks have zero sum and
direct expansion gives
\eqref{app:exact19:eq:AhelperQ}.
For the right foot of parameter
\(d\), reverse and negate the two
blocks.
In the product of an own block
with a reversed helper block only
its final two deviations survive;
their scaled product is \(-g/3\),
independently of \(k,l,m\).
Consequently for EVERY type pair,
\[
 Y_L\cdot Y_R=(g+d)/6,\qquad
          \|Y_L+Y_R\|^2=Q_L+Q_R+(g+d)/3 .
\]
We have not assumed orthogonality
of the two actual feet.

Choose \(\mu=13/100\) and form
one vector \(W=w+\mu(Y_L+Y_R)\).
At each of the two individual
anchors, the combined debit is
\[
           \mu(5g/6+1/12)
               <403/3200<79/600 .
\]
The remaining coefficient is
nonnegative before using its
original floor.
Both actual pressures and one
completed square give the following
complete three-pair table:
\[
\begin{array}{c|c|c}
 \text{types}&\|Y_L+Y_R\|^2<&\text{gain over }E>\\ \hline
 \text{finite/finite}&17/6&6487/60000\\
 \text{finite/dotted}&17729/6000&.11906\\
 \text{dotted/dotted}&9229/3000&.13000 .
\end{array}
\]
Native five belongs to the finite
column.
For example the first gain is
\(.26(.60)-.0169(17/6)\).
All three gains follow from the
full cross above and the uniform
deficit/norm bounds.
Thus
\[
            \|x\|^2>
               {1341604\over1500000}+{6487\over60000}
                     =1+{3779\over1500000}>1 .
\]
This excludes the two-weak-end case
with every internal hole retained.

\subsection{Zero or one weak end: the actual endpoints and their feet}
\label{app:exact19:Aactualends}

For either of the other reference
rows, if \(x_1\ge.55\) graft \(.55\)
into its constant first block,
preserving the block sum.
For length five the gain exceeds
\((5/4)(.55-.302)^2=.07688\);
for length six it exceeds
\((6/5)(.55-.27)^2=.09408\).
Both exceed \(1-E\) in their
respective rows.
The remaining block coordinates
stay \(>.2375\) or \(>.19\),
above the next coordinate \(<.16\);
a neighbouring hole only improves
that junction.
Reflection gives
\(x_1,-x_{20}<.55\).
The actual endpoint comparison is
\[
               x_1^2\ge(19/20)
                       \frac{\cosh L-1}{\cosh L+1}.
\]
Since \(19/60>.55^2\), both actual
endpoint cosines are less than two.

Deleting one retains \(A_{18}\).
Finite classification gives exactly
isolation; one simple-three at
\(2,3,18,19\); or one simple-four
at two or nineteen.
Two contacts make a nonpositive
cycle; exceptional finite types
cannot contain this rank-eighteen
retained path.
For \(v=g\Omega_1+\beta\Omega_j\)
the complete nonisolated data are
\[
\begin{array}{c|c|l}
 \text{white contact}&\cosh L&\text{small-own obstruction}\\ \hline
 3@2&(38g+18)/20&g=0:\ A^2<0;\quad g=1/2:\ \text{affine triangle}\\
 3@3&(38g+17)/(2\sqrt{20})&g=0:\ \text{affine }E_8\\
 3@18&(38g+2)/(2\sqrt{20})&g=0:\ A^2<0\\
 3@19&(38g+1)/20&g=0:\ A^2<0;\quad g=1/2:\ A^2=0\\
 4@2&(38g+18\sqrt2)/\sqrt{40}&\text{already }>2\text{ at }g=0\\
 4@19&(38g+\sqrt2)/\sqrt{40}&g=0:\ A^2<0 .
\end{array}
\]
For the affine \(E_8\) at old three,
the actual arms have lengths
two (old two/one), one (the cap),
and five (old four through eight).
The positive affine combination has
future time through the cap and
nonpositive omitted-wall products.
Every other positive native contact
is at least \(1/2\); in the first
row the next allowed own is at
least \(1/\sqrt2\).
Substitution makes every row
except \(3@19\) exceed two.

A remaining nonisolated left endpoint
is therefore the far-three cap
\[
 v_L=g\Omega_1+\Omega_{19}/2,\quad
 A_L^2=(19/10)g^2+g/10-21/40,\quad
                    1/\sqrt2\le g<39/38<1.1 .
\]
If both endpoints had this form,
their time squares would be
\(<(8/5)g^2,(8/5)d^2\), since
\(.3g^2+.1g-.525<0\) for \(g<1.1\).
But
\[
 v_L\cdot v_R=(gd+1/4)/10+(19/20)(g+d)
                        >(19/11)gd>(8/5)gd>A_LA_R .
\]
The actual walls are distinct:
coincidence would force \(g=d=1/2\).
Their positive mutual contradicts
acuteness.
At least one endpoint is isolated.

For an isolated endpoint,
\(v=\sigma\Omega_1\),
\(A^2=(19/10)\sigma^2-1>0\).
Own zero/three/four has nonpositive
time, six has its actual affine
\(6,3\) principal, and a parallel
own has the nonzero feasible future
null \(c+n_1\).
The complete remaining own list is
five, finite at least seven, or
dotted.
For dotted own, the actual edge
bound also gives \(19/40<u=\sigma^{-2}<1\).

Normalize each actual isolated foot
by \(\sqrt2\sigma\).
On the left write
\(Z=\zeta-P_{20}/20\), \(\sum\zeta_i=1\);
the right formula is negative reversal.
The complete profiles, pressure
\(x\cdot Z\ge t\) and square \(R=\zeta^2\)
are
\[
\begin{array}{c|c|c|c}
 \text{own}&\zeta&t&R\\ \hline
 5&(1/4,1/4,1/4,1/4)&>.34&1/4\\
 \text{finite }\ge7&(1-u/2,1-u/2,u-1)&
 \{\sqrt{1.9-u}-\sqrt{7u-4-3u^2}\}/\sqrt2&
                      3-4u+3u^2/2\\
 \text{dotted}&(1-u/2,u/2)&
 \{\sqrt{1.9-u}-\sqrt{u(1-u)}\}/\sqrt2&
                      1-u+u^2/2 .
\end{array}
\tag{A19endfeet}\label{app:exact19:eq:Aendfeet}
\]
For native five, use the finite
\(H_4\) correction
\((3/2)\sigma n_1+\sigma n_2+(\sigma/2)n_3\).
Its residual
\(\delta=1-(3/2)\sigma^2>0\)
gives the actual finite projected
point \(p+(A/\delta)(c+k)\),
of square \(-1-A^2/\delta\).
It is future and satisfies every
original wall by the selected
product calculation and positivity.
Its raw pressure is
\(H=AL/(L+\sqrt\delta)\),
\(L^2=A^2+\delta=.4\sigma^2\).
The positive-square bounds
\(A>.4934,L>.5112,\sqrt\delta<.1352\)
give \(H>.390\) and \(t>.34\).

For finite own at least seven,
use
\(c+\sigma^{-1}n_1+2(\sigma^{-1}-\sigma)n_2\).
Its cap/own products vanish, the
next product is
\((3-4\sigma^2)/(2\sigma)<0\),
and its square is
\(-(7-4\sigma^2-3/\sigma^2)<0\).
The dotted core is
\(c+\sigma^{-1}n_1\), square
\(-(1-u)\), zero cap product
and negative own product.
Both have positive future time,
positive coefficients on actual
normals, and every omitted-wall
product nonpositive.
They prove the remaining two
profiles, for this same \(x\).

The following scalar seals cover
all their parameters.
Native seven satisfies
\(.9009<\sigma<.901\), hence
\(1.23<u<1.233\),
\[
 \sqrt{1.9-u}>.8167,\quad
 \sqrt{7u-4-3u^2}<.2671,\quad
                             t>.388,\quad R<.34935 .
\]
The cosine bracket follows by signs
of \(8\sigma^3-4\sigma^2-4\sigma+1\);
it asserts no field restriction.
For finite own at least eight,
\(u<4-2\sqrt2<1.1716\),
\(\sqrt{1.9-u}>.85346\), and
\[
 7u-4-3u^2=1/12-3(u-7/6)^2<1/12 .
\]
Thus \(t>.399\), \(R\le.5\).
At \(u\ge1.1\), \(R\le.415\);
at \(u<1.1\), the radicals improve
to \(>.8944\) and \(<.265\),
giving \(t>.4448\).

For dotted own put \(q=\sqrt{1.9}\),
\(a=q-\sqrt{.9}=1/(q+\sqrt{.9})<.43\).
The lower chord and circle maximum
give
\[
 \sqrt{1.9-u}-\sqrt{u(1-u)}
       \ge q-au-\sqrt{u(1-u)}>.6185,
\]
using \(q>1.378\), \(\sqrt{.9}>.948\),
\(\sqrt{1+a^2}<1.089\).
Hence \(t>.437\).
On \(u\ge.5\) the sharper chord
\(\sqrt{1.9-u}\ge1.4178-.4692u\)
and \(\sqrt{1+.4692^2}<1.105\)
give \(t>.445\).
On \(u<.6\), use \(1.14-.5=.64\)
to get \(t>.4525\).
Also \(R<.6378125\) throughout,
and \(R\le.58\) when \(u\ge.6\).
All square-root and chord comparisons
are strict rational square checks.

On a constant end block of length
\(m\), the projection of \(Z\) is
its boundary dual divided by
\(m\sqrt2\).
Thus for \(V=Z-\Pi Z\),
\[
 w\cdot V=0,\quad \|V\|^2=R-1/m,\quad
                         \Delta_h=t-b_m .
\]
A paid coefficient \(\nu\ge0\)
costs \(\nu/(m\sqrt2)\) at that
individual boundary.
Combining the actual pressure with
the remaining original floors
BEFORE its square gives
\[
                 \|x\|^2\ge E+2\nu\Delta_h-\nu^2(R-1/m).
\tag{A19endpay}\label{app:exact19:eq:Aendpay}
\]

\subsection{No weak end: one isolated endpoint already suffices except at five}
\label{app:exact19:Azeroweak}

For zero weak ends, take \(\nu=1/2\)
in the isolated finite case.
Its price is \(<1/14<83/625\).
At native seven the gain is
\(>.086-.15/4=.0485\).
For all finite labels at least eight,
the complete subdivision is
\[
\begin{array}{c|c|c|c}
 u&t>&R-1/5<&\Delta_h-(R-1/5)/4>\\ \hline
 (1,1.1)&.4448&.3&.0678\\
 [1.1,1.14)&.4122&.215&.05645\\
 [1.14,1.1716)&.399&.1894&.04965 .
\end{array}
\]
The middle interval uses radicals
\(>.8717,<.28868\), with
\(.4122(1.4143)<.8717-.28868\).
The last uses \(>.85346,<.28868\);
the polynomial \(R\) decreases here,
with \(R(1.1)=.415,R(1.14)=.3894\).
Thus every finite isolated own
at least seven gives
\[
                         \|x\|^2>.9522798+.0485>1 .
\]

For an isolated dotted own, take
\(\nu=1/3\), price
\(<1/21<83/625\).
At \(u<.6\) the gain is
\(>(2/3)(.1505)-.438/9=31/600\).
At \(u\ge.6\) it is
\(>(2/3)(.143)-.38/9=239/4500\).
Both gains exceed \(1-.9522798\).

Only isolated native five remains.
It cannot coexist with a nonisolated
far-three opposite: its time is
\(<.497\), while the latter has
time \(<1.257\).
Their spatial product is
\(\sigma(d/10+19/20)>.95(.809)=.76855\),
whereas the time product is
\(<.497(1.257)=.624729\).
Their actual mutual would be
positive.
Thus the other endpoint is isolated
as well, and the preceding cases
force it too to have native five.

For each normalized native-five foot,
\(\Delta_h>.038\), \(\|V\|^2=.05\).
Use \(\nu=3/4\) at BOTH ends;
each price is \(<3/28<83/625\).
The two residuals have disjoint
first/last five-block support and
zero block sum, so their cross
is exactly zero.
The single vector
\(w+(3/4)(V_L+V_R)\) gives
\[
 \|x\|^2>.9522798+
       2\{(3/2)(.038)-(9/16)(.05)\}
                       =1.0100298>1 .
\]
This exhausts zero weak ends.

\subsection{One weak end: the helper and endpoint must be paid jointly}
\label{app:exact19:Aoneweak}

Orient the weak site to five.
It is far by the helper argument.
On the six/five reference its
block projection and residual are
\[
 \Pi U=(5g/6)\Omega_6+\Omega_{15}/10,\quad
 Y=U-\Pi U,\quad w\cdot Y=0,
 \qquad\Delta=H-2.288(5g/6)-2.128/10 .
\]
Relative to six/six,
the square decreases by \(1/60\)
and the deficit decreases by
\(83/3750\).
The exact finite middle-row bound
\(23165/75000\), and the stronger
native-five/dotted bounds, give
\[
                       \Delta>.2867,\qquad Q=\|Y\|^2<1.168 .
\]
The dotted upper square is
\(148/125-1/60=1751/1500<1.168\).

We must retain the full cross with
the isolated endpoint.
Write the far-helper correction
at old nineteen/eighteen as
\((l,m)=(2/3,1/3)\),
\((2(1-g^2),0)\), or \((0,0)\).
Its scaled last-five residual is
\[
             \sqrt2Y_{\rm helper}
               =(1/5,1/5,1/5-m,1/5-l+m,-4/5+l).
\]
For an isolated STRONG-side endpoint,
its residual there is \(1/5\)
minus the reversed \(\zeta\).
Direct multiplication gives
\[
\begin{array}{c|c}
 \text{isolated own}&Y\cdot V\\ \hline
 5&1/(20\sqrt2)\\
 \text{finite }\ge7&
       \{.8-u/2+m(1.5u-2)\}/\sqrt2\le(.8-u/2)/\sqrt2\\
 \text{dotted}&
       \{.8-u/2-l(1-u)-mu/2\}/\sqrt2\le(.8-u/2)/\sqrt2 .
\end{array}
\]
The finite sign uses \(u<1.233<4/3\).
For an isolated WEAK-side endpoint,
its \(\zeta\) ends before coordinate
five; multiplication with the own
six-block row printed earlier gives
\[
                             Y\cdot V=g/(3\sqrt2)<85/336 .
\]
These formulas retain every correction
coefficient, including its sign.

For paid nonnegative \(\mu,\nu\),
one vector \(W=w+\mu Y+\nu V\)
therefore yields
\[
 \|x\|^2\ge E+2\mu\Delta+2\nu\Delta_h
           -\mu^2Q-2\mu\nu(Y\cdot V)-\nu^2(R-1/m).
\tag{A19joint}\label{app:exact19:eq:Ajoint}
\]
Both individual anchor debits must
fit before either old floor is used.

If an isolated endpoint is on the
strong side, choose \(\mu=.148\).
Its own-six debit is
\[
                       \mu(5g/6)<3145/24000<79/600 .
\]
At fifteen, the helper AND endpoint
together cost
\[
       \mu/10+\nu/(5\sqrt2)
              <37/2500+4/35<83/625
\]
for all choices \(\nu\le.8\) below.
The common helper gain exceeds
\(.059279328\).
Using the full cross, the complete
table is
\[
\begin{array}{c|c|c|c|c|c}
 \text{isolated strong-side own}&\nu&
 \Delta_h>&R-1/5<&Y\cdot V<&\text{total gain over }E>\\ \hline
 5&.8&.038&.05&1/28&.07962218\\
 7&.3&.086&.15&.133&.085568928\\
 \text{finite }\ge8,\ u\ge1.1&.3&.097&.215&.179&.082234128\\
 \text{finite }\ge8,\ u<1.1&.3&.1428&.3&.215&.098867328\\
 \text{dotted},\ u\ge.6&.25&.143&.38&.358&.080537328\\
 \text{dotted},\ u<.6&.25&.1505&.438&.402&.077406328 .
\end{array}
\]
For instance the last row is exactly
\(.059279328+.5(.1505)-.0625(.438)-.074(.402)\).
All cross bounds use \(\sqrt2>7/5\)
and the stated intervals.
Every gain is \(>.0774\).

If the strong endpoint is not
isolated, the weak endpoint is
isolated by the actual endpoint
classification.
It cannot have native five:
the same time/spatial comparison
with a far-three opposite, just
proved, would force that opposite
to be isolated.
The remaining finite/dotted
weak-side cases use
\(\mu=.1,\nu=.35\).
Now BOTH feet spend the SAME
six coefficient; the full debit is
\[
 \mu(5g/6)+\nu/(6\sqrt2)
                   <17/192+1/24=25/192<79/600 .
\]
At fifteen the cost is only \(.01\).
The common helper gain is \(>.04566\)
and the full cross penalty is
\(<.07(85/336)=119/6720\).
The complete table is
\[
\begin{array}{c|c|c|c}
 \text{isolated weak-side own}&\Delta_h>&R-1/6<&
                                    \text{total gain over }E>\\ \hline
 7&.118&.183&.08813416\\
 \text{finite }\ge8&.129&1/3&.07741833\\
 \text{dotted}&.167&.472&.08703166 .
\end{array}
\]
The tight finite row is
\(.04566+.7(.129)-.1225/3-119/6720
=.077418333\ldots\).
All gains again exceed \(.0774\).
Thus in every one-weak-end case,
\[
              \|x\|^2>{3462479\over3750000}+.0774
                       =1+{2729\over3750000}>1 .
\]

\begin{proposition}\label{app:exact19:prop:wholeA19}
A whole \(A_{19}\) forest cannot occur
at the original farthest vertex of
a compact Coxeter nineteen-polytope.
\end{proposition}
\begin{proof}
The number of weak sites among the
two central end sites five/fifteen
is exactly zero, one or two.
The three preceding arguments
exclude these cases, with equality
assigned consistently to weak.
Their shared local window, endpoint
and foot lemmas precede all three
consumers and do not use any
end-case conclusion.
Every remaining internal weak
site was retained through its
actual omission, and every new
pressure was paid inside one
original-vector reference.
\end{proof}


\subsection{Whole \(B_{19}\): the real short layers before any terminal conclusion}
\label{app:exact19:Bwindow}

Use actual unit roots
\[
 n_i=(e_i-e_{i+1})/\sqrt2\ (i\le18),\qquad n_{19}=e_{19},
 \quad \omega_i=\sqrt2P_i\ (i\le18),\quad\omega_{19}=P_{19}.
\]
Now \(x_1>\cdots>x_{19}>0\).
Put \(S_i=x\cdot P_i\),
\(f_i=\sqrt i-.58/\sqrt2\),
\(d_i=\sqrt{i+1}-\sqrt i\), and
\[
                         E_K=f_5^2/5+\sum_{i=5}^{K-1}d_i^2 .
\]
Weak includes equality \(S_i\le f_i\).
The local whole-diagram lemma
\ref{app:tail:lem:whole58local} applies
literally through fifteen:
both retained old components have
rank at least four.
Its complete weak list is
\[
 v=g\omega_i+\omega_1/2,\quad
 A^2=2ig^2+2g-1/2,\quad
 \cos(\pi/5)\le g<6/5,
\tag{B19packet}\label{app:exact19:eq:Bpacket}
\]
with own five, finite at least seven,
or dotted, together with isolated
dotted five.
Pure helpers are nonadjacent.
We must extend this at the actual
\(B_3,B_2\) boundary.

At sixteen the right \(B_3\) has
a long-three extension or the
short-three \(F_4\) extension.
The latter is foreign-sealed;
own zero has the actual affine
\(F_4\) principal, whereas every
positive own has ratio
\(\ge1+1/\sqrt2>83/50\).
A full bridge is \(B_{19}\) and
directly strong.
Thus the same list holds through
sixteen, also at the old \(.61\)
threshold.
In particular some new strong
\(K=15\) or sixteen exists.
The full conditional early-reference
Lemma~\ref{app:tail:lem:whole58reference}
therefore gives a paid vector
supported through \(K\), of sum
\(f_K\), remaining individual
\(P_K\) coefficient \(>1/9\), and
effective energy
\[
               L_{15}>.94006,\qquad L_{16}>.95618,
               \quad f_{15}<3.463,\ f_{16}<3.590 .
\tag{B19early}\label{app:exact19:eq:Bearly}
\]
Effective energy means \(2D_W-\|W\|^2\),
where \(D_W\) is its actual pressure
demand; it need not equal \(\|W\|^2\).

At seventeen the right block is
\(B_2\), with simple-three at either
end.
The long-three row is strong or
affine.
For short three, own zero/four/six
have actual affine principals,
parallel has its positive future
null, and finite own at least seven
has directly strong ratio
\(\sqrt2g+1/2>83/50\).
A full bridge has residual \(1/34\)
and numerator
\(\ge g+1/34+1/(2\sqrt2)\), strong
even at own zero.
Only own three and five remain.

For own five, let \(g=\varphi/2\),
\(\alpha=\sqrt2g+1/2\).
The LEFT actual \(H_4\), using
seventeen/sixteen/fifteen, has
\[
 \delta=(5-3\varphi)/8,\quad A^2=17\alpha^2-1/2,
 \quad V_1=\cdots=V_{14}=1,\quad V^2<17,
                    \quad x\cdot V>G>4 .
\]
The point is the positive finite
projection, not its spacelike
correction; every omitted-wall
product is nonpositive.
Use \(\alpha>1.643\),
\(\sqrt\delta/\alpha<.083\)
for the last bound.
An OLD strong eleven or twelve
exists by the original bulk.
Their positive references, in both
old initial-five branches, have
\[
 q_{11}>.8472225,\ q_{12}>.8689525,\quad
 f_{11}^{\rm old}<2.886,\ f_{12}^{\rm old}<3.034,
                    \quad\text{individual endpoint price}>1/7 .
\]
These are the actual prefix
construction with its holes retained,
not an assertion that old five is
strong.
Their support lies in the flat
first fourteen.
One paid square gives
\[
 .8472225+2(4-2.886)/7-6/49>1.043,\qquad
 .8689525+2(4-3.034)/7-5/49>1.042 .
\]

For own three, the SAME cap has a
finite \(D_{19}\) replacement after
releasing the short old nineteen.
Positive inverse projection gives
an actual reciprocal short-edge
endpoint, with
\[
 \cosh L_{19}=17/2+19\sqrt2/4>15,\qquad
               S_{19}^2>133/8>(4.07)^2 .
\]
Use \eqref{app:exact19:eq:Bearly},
and the pressure \(Z=P_{19}\).
The legal \(1/9\) endpoint payment
gives
\[
 .94006+2(4.07-3.463)/9-4/81>1.025,\qquad
 .95618+2(4.07-3.590)/9-3/81>1.025 .
\]
Thus coverage reaches seventeen.

Nonadjacency is also actual at this
boundary: replacing two adjacent
helper caps retains the long
\(A_{16}\) with its two simple
leaves, a positive \(D_{18}\), and
the independent old short \(A_1\).
The positive two-cap inverse and
all-wall projection therefore apply.
No long nineteenth old normal is
invented.
Every weak site five through
seventeen is now
\eqref{app:exact19:eq:Bpacket},
except isolated dotted five.

\subsection{An isolated long-end foot is a local geometric kernel}
\label{app:exact19:longfeet}

We record the point formulas once.
They apply to any actual isolated
long-end cap with initial long roots
one/two/three as above, in \(B\) or
\(D\), whenever the listed own occurs.
Write
\[
                  c_L=Ap-\gamma\omega_1,\qquad A^2=2\gamma^2-1 .
\]
For each remaining own type the
actual normalized pressure is
\(x\cdot Z\ge h\), \(\sum Z_j=1\):
\[
\begin{array}{c|c|c|c}
 \text{own}&Z&h&R=Z^2\\ \hline
 5&P_4/4&\sqrt{(\varphi-1)/4}&1/4\\
 \text{finite }\ge7&
 (1-u/2)(e_1+e_2)+(u-1)e_3&
 \{\sqrt{2-u}-\sqrt{7u-4-3u^2}\}/\sqrt2&
                   3-4u+3u^2/2\\
 \text{dotted}&(1-u/2)e_1+(u/2)e_2&
 \{\sqrt{2-u}-\sqrt{u(1-u)}\}/\sqrt2&
                   1-u+u^2/2 .
\end{array}
\tag{B19headfeet}\label{app:exact19:eq:longfeet}
\]
Here \(u=\gamma^{-2}\), in
\((1,\cos(\pi/7)^{-2}]\) for finite
own, and \((0,1)\) for dotted.
All coordinates of \(Z\) are
nonnegative.

At native five, the actual \(H_4\)
correction is
\(\gamma(3n_1/2+n_2+n_3/2)\),
with residual
\(\delta=1-3\gamma^2/2>0\).
The future finite projected point
\(p+(A/\delta)(c_L+k)\) has square
\(-1-A^2/\delta\) and every original
wall product nonpositive.
Its displacement is
\((\gamma/4)\omega_4\), and its
pressure is \(H=AB/(B+\sqrt\delta)\),
\(B=\gamma/\sqrt2\).
The real native-five identity
\(4\gamma^2=2\gamma+1\) gives
\(\sqrt\delta=(1-\gamma)/\sqrt2\),
hence \(H=\gamma A\) and the first
row exactly.
In particular \(h>.393075\),
by positive rational squaring.

For finite own at least seven,
the actual negative core is
\[
 c_L+\gamma^{-1}n_1+2(\gamma^{-1}-\gamma)n_2 .
\]
Its cap/own products vanish, the
next product is
\((3-4\gamma^2)/(2\gamma)<0\),
and its square is
\(-(7-4\gamma^2-3/\gamma^2)<0\).
For dotted own use
\(c_L+\gamma^{-1}n_1\), square
\(-(1-u)\), zero cap product and
negative own product.
Both have positive future time
and positive coefficients on actual
normals.
Every omitted original-wall product
is nonpositive termwise.
Their spatial displacements, divided
by \(\sqrt2\gamma\), give the other
two rows on the SAME original \(x\).
This proves the kernel without an
\(I\) spectator or an arithmetic
hypothesis.

For every finite own at least seven
and every dotted own,
\[
                              .8h-.16R>.2893 .
\tag{longF}\label{app:exact19:eq:longF}
\]
At native seven,
\(.9009<\gamma<.901\) gives
\(A>.7894\),
\(\sqrt{7-4\gamma^2-3/\gamma^2}<.2395\),
\(\sqrt2\gamma<1.2743\);
thus \(h>.4315\).
Also \(1.2318<u<1.233\) and
\(R<.349\), giving \(>.28936\).

For finite labels at least eight,
\(1<u<1.172\).
Put \(a=.8/\sqrt2\),
\(D=(u-1)(4-3u)\).
The derivative of the left side is
\[
 a\left\{-\frac1{2\sqrt{2-u}}-\frac{7-6u}{2\sqrt D}\right\}
                         +.64-.48u .
\]
On \((1,7/6]\), the first term is
\(<-.28\), the second is nonpositive,
and the last is at most \(.16\).
On \([7/6,1.172)\), \(\sqrt D>.284\),
the positive second term is
\(<.032\), and the other terms
are \(<-.28+.08\).
Thus the scalar decreases with
\(u\); its minimum is at native
eight.
There the radical bounds are
\(>.910,<.289\), and \(R<.373\).
Hence \(h>.4388\) and the scalar
is \(>.29136\).

For dotted own,
\(.5656<a<.5657\).
On \(u\le.5\), use
\(\sqrt{2-u}\ge1.4142-.379u\),
\(R\le1-.75u\), giving
\[
 .8h-.16R>.63987152-.0944u-.5657\sqrt{u(1-u)}
                                      >.30587152 .
\]
The circle bound uses
\(\sqrt{.0944^2+.5657^2}<.5736\).
On \(u\ge.5\), use
\(\sqrt{2-u}\ge1.4494-.4494u\),
\(R\le.75-.25u\), giving
\[
 .8h-.16R>.69978064-.2142u-.5657\sqrt{u(1-u)}
                                      >.29023064 ,
\]
with \(\sqrt{.2142^2+.5657^2}<.6049\).
The chords lie below the concave
root by their endpoint square
checks.
This proves the entire continuous
dotted range.

\subsection{The \(B_{19}\) long head and its paid gains}
\label{app:exact19:Bhead}

The actual long head is isolated.
Indeed deleting one retains
\(B_{18}\), whose complete finite
one-node list is isolation or
simple-three at its long leaf two.
A second contact makes a cycle;
another four or an additional fork
is not finite; exceptional finite
ranks are too small.
The nonisolated cap would have
\[
 v=\gamma\omega_1+\omega_2/2,\quad
 A^2=2\gamma^2+2\gamma,\quad\cosh L_1=2\gamma+1 .
\]
Own zero has zero time and own
three is the actual affine
triangle, so \(\gamma\ge1/\sqrt2\)
and \(x_1>\sqrt{\sqrt2-1}>.643\).

The OLD actual bulk through fifteen
has some strong fourteen/fifteen,
with a paid reference of energy
\(>.90715\), every weak hole omitted.
If old five is strong, graft
\(.643\) into its first-five block
of mean \(<.361\).
The other four coordinates remain
\(>.2904\), above the outgoing
slope \(<.214\), and the gain is
\((5/4)(.643-.361)^2=.099405\).
If old five is weak, six is strong.
Replace the first-six block by
\(.643\) and five equal slopes
of sum \(f_6-.643\).
Its energy is
\[
 .643^2+(2.01815-.643)^2/5=.7916565045 ,
\]
whereas the old first-six energy
is \(<.6971\).
The gain \(>.09455\), with the
same retained tail, again gives
energy above one.
Both references are ordered and
paid by the original \(x\).
Thus the nonisolated head is
impossible.

For the isolated cap,
\(A^2=2\gamma^2-1>0\) excludes
own zero/three/four.
Own six has the affine \(6,3\)
principal on the actual first two
roots; parallel has its feasible
positive future null.
The remaining owns are exactly
native five, finite at least seven,
or dotted.
Their actual feet are
\eqref{app:exact19:eq:longfeet}.

A native-five head forces true
new strong five/six/seven.
For a helper at \(i\le7\),
\[
 c_L\cdot c_i=\gamma_0(2g+1)-A_0A_i,\qquad
       \gamma_0=\varphi/2,\quad A_0^2=(\varphi-1)/2 .
\]
Since \(\gamma_0^2/A_0^2=\varphi^3/2>21/10\),
\[
 A_i^2-\frac{\gamma_0^2}{A_0^2}(2g+1)^2
                   <(28g^2-32g-13)/5<0
                  \qquad(.8<g<1.2).
\]
The convex quadratic is negative
at both endpoints.
The actual mutual would be
positive.
For the remaining weak isolated
dotted five,
\[
 c_L\cdot c_5=\varphi g-A_0\sqrt{10g^2-1}.
\]
Its sign is strictly negative,
since the squared difference in
the opposite order is
\((4\varphi-6)g^2-(\varphi-1)/2
\ge(7\varphi-11)/2>0\).
Weakness gives \(g<1/.58<2\);
\(\varphi-\sqrt{5(\varphi-1)}>-3/20\)
then bounds the mutual in
\((-3/10,0)\), a forbidden native
interval.
This proves the stated three
strong sites without a later
terminal floor.

For any strong-five \(.58\)
reference with first-five mean
\(b=f_5/5<.3652\) and individual
\(P_5\) coefficient \(>.151\),
a head correction
\(\mu(Z-P_5/5)\) costs \(\mu/5\)
before its floor is used.
It is orthogonal to that constant
block and has square \(R-1/5\).
Its effective gain is
\[
                             G_L=2\mu(h-b)-\mu^2(R-1/5).
\]
The complete paid gains are
\[
\begin{array}{c|c|c}
 \text{head}&\mu&G_L>\\ \hline
 5&.56&.015456\\
 \text{finite }\ge7&.4&.02914\\
 \text{dotted}&\sqrt2/5&.03384 .
\end{array}
\tag{B19headgains}\label{app:exact19:eq:Bheadgains}
\]
The first uses \(h>.393\);
the second is
\eqref{app:exact19:eq:longF}
minus \(.8b\) plus \(.032\).
For the dotted row, expand its
gain as
\[
 .4\{\sqrt{2-u}-\sqrt{u(1-u)}-\sqrt2b\}
                                -.04(1.6-2u+u^2),
\]
where \(\sqrt2b<.5166\).
For \(u\le.5\), its lower bound is
\(.29504-.0916u-.4\sqrt{u(1-u)}>.04404\),
using \(\sqrt{.0916^2+.4^2}<.4104\).
For \(u\ge.5\), it is
\(.32912-.15976u-.4\sqrt{u(1-u)}>.03384\),
using \(\sqrt{.15976^2+.4^2}<.4308\).
The same two concave root chords
prove these inequalities.
All three individual prices fit.

\subsection{Curved early references with the whole initial-five debt}
\label{app:exact19:Bcurvedearly}

Some subsequent pressures need a
stronger early reference than
\eqref{app:exact19:eq:Bearly}.
We give it by the actual initial
five alternatives, retaining all
holes and full cross terms.
Each resulting vector is supported
through \(K=15\) or sixteen,
has sum \(f_K\), and keeps endpoint
coefficient \(>1/9\).

First record a useful exact scalar:
\[
                              E_{17}>.9726645 .
\tag{B19E17}\label{app:exact19:eq:BE17}
\]
Its first block is
\(f_5^2/5=1.03364-.116\sqrt{10}>.6668157\),
using \(\sqrt{10}<3.162278\).
For \(a=4i+2\),
\(u=d_i^2\) satisfies \(u(a-u)=1\).
Three positive reciprocal
substitutions give
\(u>(a^2-1)/(a(a^2-2))\).
At denominator \(10^7\), the full
list of lower numerators and
cleared positive margins is
\[
\begin{array}{c|r|r@{\qquad}c|r|r}
 i&L_i&\text{margin}&i&L_i&\text{margin}\\ \hline
5&455488&5248&11&217494&13464\\
6&385186&536&12&200080&8000\\
7&333704&14240&13&185248&115712\\
8&294372&20208&14&172465&14860\\
9&263340&21360&15&161332&72272\\
10&238230&27080&16&151549&273164 .
\end{array}
\]
The margin is
\(10^7(a^2-1)-L_i a(a^2-2)\).
The numerators sum to \(3058488\),
proving \eqref{app:exact19:eq:BE17}.
This is a finite rational hand
certificate, not a script-output
premise.
Also \(d_{17}^2>1/70\), giving
\(E_{18}>.98659\).

If five is strong, its curved
reference loses \(<1/1024\) over
all nonadjacent internal holes.
Adding the actual head pressure
inside that reference gives
\[
                   L_{15}>.955516,\qquad L_{16}>.971636 .
\]
These conservative values use the
smallest gain in
\eqref{app:exact19:eq:Bheadgains};
the appropriate finite/dotted head
is stronger.

If five is a weak helper, the
actual packet gives \(x_1>.5455\)
and six is strong.
Use the first-six reference through
these actual heights, followed by
the ENTIRE curved retained window.
Its first-six energy is \(>.74385\),
rather than the ordinary \(<.7126\).
Since \(d_{15}^2<.016134\),
\(d_{16}^2<.015155\),
\eqref{app:exact19:eq:BE17} gives
\(E_{15}>.9413755,E_{16}>.9575095\).
After all holes,
\[
                          L_{15}>.9716,\qquad L_{16}>.9877 .
\]
No long-head energy was added to
this initial-coordinate gain.

If five is isolated dotted,
the native-five head has already
been excluded by its compatibility.
For the other head types keep its
actual foot \(x\cdot Z\ge h\)
and the initial-five foot
\[
 Z_5=P_5+(v/2)(e_6-e_5),\quad
                x\cdot Z_5>T_5>1.8185,\quad .4<v<1 .
\]
This is the actual dotted-five
negative core; the weak condition
forces \(v>.4\) by the full
continuous loss bound in the
shared early-reference lemma.
Choose a genuinely strong
\(J=6\) or seven.
The reference is constant
\(m=f_J/J\) on the first \(J\)
coordinates and follows every
retained curved knot through \(K\).
The complete conservative data are
\[
\begin{array}{c|c|c|c|c|c}
 J&m<&w_{15}^2>&w_{16}^2>&
 \text{individual }P_J\text{ coefficient}>&Q_5<\\ \hline
 6&.3399&.9217&.9378&.1435&77/150\\
 7&.31938&.9041&.9202&.1366&194/175 .
\end{array}
\]
For the energies, \(f_6^2/6>.69317\),
\(f_7^2/7>.71400\).
The bound \(d_i^2>1/(4i+2)\) gives
at denominator \(10^6\) the lower
numerators
\[
38461,33333,29411,26315,23809,21739,
20000,18518,17241,16129
\]
for \(i=6,\ldots,15\).
Those from six through fourteen sum to \(228827\),
and from seven to fourteen to
\(190366\); \(d_{15}^2>.016129\).
The maximal omission losses are
\(<.000267\) after six and
\(<.000181\) after seven, by the
decreasing sums \(1/[32(j-1)^3]\)
over their earliest permitted
nonadjacent sets.
These prove every energy entry.
The \(P_J\) price is \(m-d_J\);
the end-\(K\) price stays \(>1/9\).
Finally
\[
 Q_5=\|Z_5-(5/J)P_J\|^2
          =5-25/J-v+v^2/2<4.68-25/J .
\]

Let \(Y=Z-P_J/J\),
\(Y_5=Z_5-(5/J)P_J\).
Their first-\(J\) sums vanish,
so they are orthogonal to \(w\).
Because \(Z\) is supported on the
first three coordinates and
\(Z_5=1\) there,
\[
             Y^2=R-1/J,\qquad Y\cdot Y_5=1-5/J .
\]
The one vector \(w+.4Y+.09Y_5\)
spends the COMBINED
\(P_J\) price \(.85/J\), which
fits both rows.
Its gain is
\[
 .8h-.16R-.8m+.16/J
 +.18(T_5-5m)-.0081Q_5-.072(1-5/J).
\]
The nonzero last cross is retained.
Using \eqref{app:exact19:eq:longF},
the gain is \(>.04930\) at six
and \(>.06699\) at seven.
Thus both initial branches give
\[
                         L_{15}>.9710,\qquad L_{16}>.9871 .
\]
These are effective energies of
the same paid vector, not new
independent norm summands.

The identical curved construction
at a genuinely strong seventeen
has initial reference squares
\(>.9530,.9354\) in the last case,
so its gains give totals
\(>1.00230,1.00239\).
For helper five it gives
\(>.97231+.03125-1/1024>1.0025\).
At a later actual strong eighteen,
the extra endpoint increment
dominates the possible additional
hole, so the conclusions remain
valid.
These are conditional uses of
an ACTUAL strong endpoint;
none is used to assume that
endpoint beforehand.

% Complete B19 terminal classification/zero-own/survivor proofs follow.


